\chapter{Geometric Extraction of the Local Navigation Reference}
\label{cap:geometric_reference}

This chapter describes the geometric extractor implemented in
\texttt{pf4.py} (referred to as Method~4 in the experimental evaluation).
It converts a segmented corridor mask into a local navigation reference
by extracting intervals in horizontal strips, associating their left and
right boundaries from bottom to top, and fitting each edge independently
with a robust linear estimator. The centerline reference is the analytic
mean of the two fitted edge lines.


\section{Input representation and edge-line models}
\label{sec:c6_input}

Both functions expect a two-dimensional label mask $M(y,x)$ of height $H$
and width $W$. Image coordinates increase to the right in $x$ and downward
in $y$. The argument \texttt{path\_class} identifies the class from which
corridor intervals are extracted; its default is 1 in both files. A
binary input can therefore use 1 for the target region and 0 elsewhere.
For a multiclass segmentation mask, this argument must match the actual
label encoding, or the mask must first be converted appropriately. The
extractors do not infer the class mapping or use a separate plant class.
A target-region label is not, by itself, a guarantee of traversable space.

The left and right boundaries are represented by
\begin{equation}
    x_L(y)=m_Ly+q_L, \qquad x_R(y)=m_Ry+q_R.
    \label{eq:c6_edges}
\end{equation}
Treating $x$ as a function of $y$ permits vertical image lines without an
infinite slope. The center reference is
\begin{equation}
    x_C(y)=\frac{x_L(y)+x_R(y)}{2}=m_Cy+q_C,
    \qquad m_C=\frac{m_L+m_R}{2},\quad q_C=\frac{q_L+q_R}{2}.
    \label{eq:c6_line}
\end{equation}
Thus, both methods assume locally straight image-space boundaries.
The models do not reconstruct terrain or estimate a metric trajectory.
Their residuals and width estimates are in pixels, not lateral distances
in meters. The two edges need not be parallel: their separation can
change linearly with image height.

At each image height, the reference lies halfway between the two
predicted boundaries. This construction is not generally equivalent to
fitting the observed interval midpoints: the two edge fits can exclude
different border-touching samples and assign different residual weights.
Keeping the boundaries separate allows each visible side to contribute
its own observations, although both sides must have enough samples.

The shared processing stages are:
\begin{enumerate}
    \item select a lower region of interest and clean the target-class mask;
    \item extract and merge projected intervals in horizontal strips;
    \item select a seed using the method-specific support and priority rules;
    \item associate interval boundaries upward through the image;
    \item fit the non-clipped left and right boundaries independently;
    \item calculate quality indicators and rasterize the center reference,
    with method-specific validation and output rules.
\end{enumerate}

\section{Region of interest and mask preparation}
\label{sec:c6_roi}

The ROI begins at $y_0=\lfloor\rho H\rfloor$ and extends to the bottom
of the image. The default starting ratio is $\rho=0.30$ for Method~4.
The binary target mask is
\begin{equation}
    B(y,x)=\mathbf{1}\{M(y_0+y,x)=\texttt{path\_class}\},
    \qquad 0\leq y<H_R=H-y_0.
\end{equation}
Opening with a $5\times5$ elliptical kernel is applied first,
then closing with a $9\times9$ rectangular kernel. Connected components
with fewer than 500 pixels are removed; all sufficiently large components
are retained, not only the largest one.

Let $\widetilde B$ denote the cleaned binary ROI mask. It supplies both
interval extraction and local support calculations. These support
calculations do not use the original multiclass mask to compute separate
soil and plant ratios.

Morphology and interval merging can bridge interruptions or change region
connectivity. Consequently, support in $\widetilde B$ measures agreement
with the processed mask, not independently verified scene geometry.
Kernel sizes, component areas, and distance thresholds are in input-mask
pixels and depend on image resolution and camera configuration.

\section{Horizontal strips and candidate intervals}
\label{sec:c6_candidates}

\subsection{Strip construction and column projection}
\label{subsec:c6_strips}

Both methods use $N=24$ strips by default, indexed from bottom to top.
Their height is $h=\max(1,\lfloor H_R/N\rfloor)$. For strip index $s$,
the ROI-relative bounds are
\begin{equation}
    y_s^- = \max(0,H_R-(s+1)h),\qquad
    y_s^+ = \min(H_R,H_R-sh),
\end{equation}
with the lower image row bound $y_s^+$ excluded. Empty strips are skipped.
The representative full-image coordinate is
$y_s=y_0+(y_s^-+y_s^+)/2$, following the code's bound-based convention.
Integer division may leave a remainder above the constructed strips.

For example, with $H_R=157$ and $N=24$, the strips cover the lowest
$24\times6=144$ ROI rows. The upper 13 ROI rows do not supply candidates;
the ROI and the region actually sampled by the strips are therefore not
identical.

Within each strip, the cleaned mask is summed vertically. A column is
active whenever its sum exceeds zero, so even one target pixel activates
that column. Contiguous active columns form raw intervals. Neighboring
intervals are merged when the number of inactive columns between them
is at most \texttt{max\_interval\_gap}, whose default is 30 pixels.
For a merged interval $[\ell_i,r_i]$, the width is
\begin{equation}
    d_i=r_i-\ell_i+1.
    \label{eq:c6_interval}
\end{equation}
Intervals narrower than 8 pixels are discarded. There is no fixed maximum
interval width in either new method. Instead, an interval touching both
image margins is rejected if the mean of the entire cleaned binary strip
is below 0.70. With \texttt{border\_margin}=2, the flags are

\placeholderfigure{fig:geometric_strips_projection}%
{Horizontal strip decomposition and column-wise projection for path interval detection}%
{Schema illustrativo della decomposizione della maschera binaria di navigazione in $N=24$ strisce orizzontali e relativa proiezione delle colonne attive per l'identificazione degli intervalli di corridoio.}%
{5cm}
\begin{equation}
    t_{L,i}=\mathbf{1}\{\ell_i\leq2\},\qquad
    t_{R,i}=\mathbf{1}\{r_i\geq W-1-2\}.
\end{equation}
These flags indicate possible truncation by the image boundary, rather
than verified physical corridor edges.

Column projection merges evidence across different vertical locations;
a merged interval can also contain horizontal gaps. Neither its width
nor its endpoints prove that its rectangular interior is entirely target
class. This matters when interpreting intervals around occlusions or
fragmented segmentation.

\subsection{Local support around interval boundaries}
\label{subsec:c6_support}

Method~4 computes $u_{L,i}$ and $u_{R,i}$ in separate bands centered on
$\ell_i$ and $r_i$. The effective half-width is
$\max(2,\texttt{local\_band\_half\_width})$, equal to 5 by default.
These means quantify target-pixel occupancy around each boundary; they
are not boundary-detection probabilities. Because a band straddles a
boundary, it need not be entirely target class even for a clean interval.

Edge support ranks seed candidates and contributes to the final quality
score. It is not an acceptance threshold and does not weight the line fit.
In particular, \texttt{min\_edge\_support}=0.20 is present in the
function signature but unused in the active implementation.

\section{Initialization and bottom-up association}
\label{sec:c6_selection}

\subsection{Edge-support seed selection and gap-limited association}
\label{subsec:c6_seed4}

Method~4 pools all candidates in the first six constructed strips from
the bottom and selects the seed by maximizing $\min(u_{L,i},u_{R,i})$,
then interval width. No minimum-support threshold is applied, so even a
weakly supported candidate can initialize tracking.

The boundary displacement $\delta_i$ and association cost are
\begin{equation}
    D_4=\max(28,0.5d_{\mathrm{prev}}),\qquad
    J_{4,i}=\delta_i+0.4|d_i-d_{\mathrm{prev}}|.
    \label{eq:c6_cost4}
\end{equation}
Only the displacement limit determines admissibility; edge support does
not enter this cost or an additional tracking gate.

A counter records consecutive strips without a compatible candidate and
resets after each successful association. With
\texttt{max\_gap\_strips}=3, tracking stops on the fourth consecutive
miss. Stopping does not directly return \texttt{lost}: the accumulated
points proceed to fitting if at least three are available. The maximum
miss count, including the miss that triggers stopping, enters the quality
score. The previous interval remains unchanged during gaps.

The stopping rule and gap penalty have different effects: three misses
are tolerated for continued association, but already reduce the gap
score to zero. A fourth miss stops the search without automatically
rejecting a reference supported by the observations collected earlier.

Method~4 performs greedy spatial association within one frame, not
temporal tracking. Each accepted interval contributes two boundary
observations, and at least three accepted intervals are required before
the edge-fitting stage.

\section{Separate robust fitting of the corridor edges}
\label{sec:c6_fitting}

\subsection{Excluding image-boundary observations}
\label{subsec:c6_edge_samples}

The left fit uses only accepted intervals not flagged as touching the
left image margin. The right fit independently excludes intervals
flagged as touching the right margin. At least two observations per edge
must remain. The two fits can therefore use different strips and span
different vertical ranges. A clipped boundary is not replaced by a
one-sided centerline estimate: if either edge lacks sufficient samples,
both implementations return an empty output with \texttt{lost} status.

The retained samples are interval endpoints from the cleaned column
projection. They are not points obtained by a dedicated edge detector.
Excluding border-touching samples reduces the influence of image
truncation, but does not establish the correctness of the remaining
boundaries.

\subsection{Residual-based reweighting}
\label{subsec:c6_robust_fit}

Both methods start each edge fit with uniform weights. There is no
semantic, interval-width, or vertical-position base weighting. For an
edge sample $(x_i,y_i)$, four reweighting iterations are performed by
default, with scale $\sigma=18$ pixels:
\begin{equation}
    e_i^{(t)}=x_i-m^{(t)}y_i-q^{(t)},\qquad
    w_i^{(t+1)}=\frac{1}{1+(e_i^{(t)}/\sigma)^2}.
    \label{eq:c6_reweight}
\end{equation}
Each iteration calls \texttt{numpy.polyfit} with degree one and updated
weights. Its weights multiply unsquared residuals
\cite{numpyPolyfitDocumentation}; the fitted objective at that iteration is
\begin{equation}
    (m,q)=\mathop{\arg\min}_{m,q}
        \sum_i \left[w_i(x_i-my_i-q)\right]^2.
    \label{eq:c6_objective}
\end{equation}
Thus, the effective coefficient of each squared residual is $w_i^2$,
not $w_i$. The implementation should not be interpreted as a different
standard robust-loss estimator solely from the form of its weight rule.
Weights are recomputed from the current model rather than accumulated
across iterations; no convergence-based stopping criterion is used.

\placeholderfigure{fig:robust_edge_fitting}%
{Robust line fitting of corridor boundaries and centerline construction}%
{Diagramma del processo di fitting robusto separato dei bordi sinistro ($x_L$) e destro ($x_R$) tramite re-weighting iterativo dei residui con fattore lorentziano ($\sigma=18$~px) e calcolo dell'asse centrale di guida $x_C(y)$.}%
{5cm}

The final residual for each retained sample is its absolute horizontal
deviation from the corresponding fitted edge. Define $\bar e$ as the
mean over the concatenated left and right residual arrays, and
\begin{equation}
    F=\operatorname{clip}(1-\bar e/\sigma,0,1).
    \label{eq:c6_fit_score}
\end{equation}
Edges with more retained samples contribute more terms to this mean.
Reweighting reduces the influence of inconsistent points but cannot
correct an erroneous association chain that supports a coherent wrong
line. It does not impose parallel edges or estimate genuine path curvature.

\section{Quality indicators, validation, and image-space output}
\label{sec:c6_output}

\subsection{Method 4: geometric checks and a composite score}
\label{subsec:c6_output4}

Method~4 validates the fitted boundaries before constructing the output.
At 16 evenly spaced heights covering the combined range of retained
edge samples, it evaluates
\begin{equation}
    D(y)=x_R(y)-x_L(y).
\end{equation}
All sampled widths must be strictly positive. The median of these
sampled widths is returned as \texttt{width}. This check concerns the
sampled observation range, not any extrapolation below it.

A separate indicator uses all accepted intervals, including those with
border-touching endpoints. For $K$ intervals in association order, let
$z_j=r_j-\ell_j$; unlike extraction width $d_j$, this quantity omits the
inclusive-endpoint $+1$. The width-irregularity indicator is
\begin{equation}
    E_D=\frac{Q_{0.90}\left(\{|z_{j+2}-2z_{j+1}+z_j|\}_{j=1}^{K-2}\right)}
                   {\max(\operatorname{median}_j|z_j|,1)},
    \qquad A=\operatorname{clip}(1-E_D/0.30,0,1),
    \label{eq:c6_width_score}
\end{equation}
where $Q_{0.90}$ is the 90th percentile. These are second differences in
sample order, not derivatives normalized by vertical distance. If strips
are skipped, sample spacing is unequal; even a width varying linearly
with image height need not yield zero second differences. The indicator
therefore measures an implementation-specific irregularity, not a
perspective-invariant physical curvature.

Let $S$ be the mean of the left support values used by the left fit and
the right support values used by the right fit, pooled together. With
maximum consecutive miss count $g_{\max}$ and allowed gap $g_a=3$,
\begin{equation}
    G=\operatorname{clip}(1-g_{\max}/\max(g_a,1),0,1),\qquad
    C_4=0.25T+0.20S+0.25F+0.20A+0.10G.
    \label{eq:c6_confidence4}
\end{equation}
The result is \texttt{lost} if the positive-width check fails,
$\bar e>2\sigma$, or $C_4<0.38$. Otherwise, it is \texttt{degraded}
for $C_4<0.60$ and \texttt{ok} for $C_4\geq0.60$.

A result marked \texttt{lost} returns a zero output matrix before
rasterization, with no centerline model in the post-fit rejection branch.
For an accepted result, the mean edge line is rasterized from the top of
the highest accepted strip to the bottom of the image, then dilated with
a $3\times3$ elliptical kernel. This can extrapolate below the seed. The
edge-order check does not cover the entire rendered range when that range
exceeds the sampled observation range.

Diagnostics report the edge and centerline models, mean residual,
support, tracking fraction, median fitted width, width irregularity,
maximum gap, row bounds, and cleaned ROI mask. Earlier failures return
\texttt{lost} with reason strings such as \texttt{no\_seed\_candidate},
\texttt{tracking\_failed}, or \texttt{insufficient\_edge\_points}.

\subsection{Interpretation of status and confidence}
\label{subsec:c6_failures}

The confidence score $C_4$ is a heuristic combining processed-mask edge
support with geometric and gap indicators. It is not a calibrated
probability. A downstream controller must inspect the returned status;
neither function specifies a robot stopping or recovery policy.
Malformed inputs, unsuitable parameter values, or numerical exceptions
are not comprehensively handled by the explicit failure branches.

\section{Default configuration}
\label{sec:c6_discussion}

Table~\ref{tab:c6_defaults} lists the default parameter values for Method~4.
These are the defaults of \texttt{pf4.py}, not a statement that a
controller uses them or that they are experimentally optimal. Overrides
and mask encodings must be recorded for each experiment.

\begin{table}[htbp]
    \centering
    \small
    \caption{Default configuration of Method~4 (\texttt{pf4.py}).
    Distances and areas are expressed in input-mask pixels.}
    \label{tab:c6_defaults}
    \begin{tabularx}{\textwidth}{@{}X l@{}}
        \toprule
        Setting & Value \\
        \midrule
        Target-class identifier & 1 \\
        ROI starting ratio $\rho$ & 0.30 \\
        Number of strips $N$ & 24 \\
        Opening / closing kernels & $5\times5$ / $9\times9$ \\
        Minimum component area & 500 px \\
        Minimum interval width & 8 px \\
        Maximum merged gap & 30 px \\
        Image-border margin & 2 px \\
        Full-strip occupancy gate & 0.70 \\
        Local support-band half-width & 5 px \\
        Seed search window & 6 strips \\
        Jump-limit lower bound & 28 px \\
        Width-change cost coefficient & 0.4 \\
        Consecutive misses allowed & 3 \\
        Residual scale $\sigma$ / iterations & 18 px / 4 \\
        Width-error score scale & 0.30 \\
        Lost / acceptable score thresholds & 0.38 / 0.60 \\
        Output dilation kernel & $3\times3$ \\
        \bottomrule
    \end{tabularx}
\end{table}

The two-line representation retains boundary information separately,
but requires sufficient non-clipped samples on both sides. Independent
edge fits may be supported over different height ranges, and their mean
can include extrapolation of one or both boundaries. A wide or
well-supported interval is not necessarily the intended row.

\section{Effect of strip height and number}
\label{sec:c6_strip_height}

The strip count controls vertical sampling and the maximum number of
interval observations. More strips provide thinner vertical regions and
potentially more points per boundary fit. Taller strips aggregate more
mask pixels but mix evidence from different image heights, and column
projection can conceal interruptions. Increased strip height does not
guarantee more reliable boundary estimates or a numerically better fit.

For a $224\times224$ input with $\rho=0.30$, Method~4 has $y_0=67$ and
$H_R=157$. Table~\ref{tab:c6_strip_tradeoff} shows strip heights for
different strip counts using these dimensions. At $N=24$, the strip
height is 6 pixels. A fit can use fewer than $N$ points because of seed
location, missed strips, early stopping, or border-point exclusion;
the two edges can have different sample counts.

\begin{table}[htbp]
    \centering
    \small
    \caption{Strip heights for an illustrative $224\times224$ mask with
    $\rho=0.30$ ($H_R=157$). The point count is an upper bound per edge,
    not a guaranteed number of fitted samples.}
    \label{tab:c6_strip_tradeoff}
    \begin{tabular}{@{}c c c@{}}
        \toprule
        Strips $N$ & Strip height (px) & Max. edge samples \\
        \midrule
        24 & 6 & 24 \\
        20 & 7 & 20 \\
        16 & 9 & 16 \\
        12 & 13 & 12 \\
        8 & 19 & 8 \\
        \bottomrule
    \end{tabular}
\end{table}

Changing $N$ also changes the vertical extent of a fixed six-strip seed
window and the extent of a three-strip tolerated gap. These settings
should therefore be considered together. The width-irregularity indicator
depends on the sampling sequence, while tracking fractions use the number
of constructed strips as their denominator. Selecting $N$ requires
evaluating boundary and centerline errors, accepted-point counts,
rejection rates, and processing time under the intended image resolution
and field conditions. No optimal strip count is inferred from the
implementation alone.

\section{Baseline algorithms for comparative evaluation}
\label{sec:c6_baseline_algorithms}

To quantify the performance of decoupled dual-edge regression against alternative geometric strategies, three baseline extractors are integrated into the evaluation pipeline: a center-point strip tracker, a column-histogram extractor, and a RANSAC-based dual-edge extractor.

\subsection{Center-point strip tracking baseline}
\label{subsec:c6_center_baseline}

The center-point strip tracker (\texttt{pf\_centri.py}) shares the pre-processing stage, region of interest ($y_0 = \lfloor 0.30 H \rfloor$), and horizontal strip decomposition ($N = 24$) of Method~4. In each strip, horizontal intervals are extracted from the binary corridor mask. Instead of retaining the left and right boundary coordinates independently, the algorithm computes the midpoint and width of each candidate interval:
\begin{equation}
    x_C = \frac{x_L + x_R}{2}, \qquad w = x_R - x_L + 1.
    \label{eq:c6_center_def}
\end{equation}

Tracking starts from the candidate interval with the highest pixel density within the bottom six strips. Successive intervals are associated bottom-up by minimizing a transition cost combining lateral center displacement and corridor width variation:
\begin{equation}
    J_{\text{center}} = |x_C^{(k+1)} - x_C^{(k)}| + 0.4\, |w^{(k+1)} - w^{(k)}|,
    \label{eq:c6_center_cost}
\end{equation}
subject to a maximum allowable jump $|x_C^{(k+1)} - x_C^{(k)}| \le \max(28, 0.5\, w^{(k)})$.

Once tracking reaches the upper strip or exceeds the maximum gap of three consecutive empty strips, a single centerline model $x_C(y) = m_C y + q_C$ is fitted to the collected centers $(y_k, x_{C,k})$ using IRLS with Cauchy weighting across four iterations. Because interval midpoints are formed prior to fitting, asymmetric border disturbances shift $x_{C,k}$ directly, transferring localized boundary noise into the fitted trajectory.

\subsection{Column-histogram baseline}
\label{subsec:c6_histogram_baseline}

The column-histogram baseline (\texttt{pf\_histogram.py}) reproduces the SegMin methodology developed by Navone et al.~\cite{navone2025gpsfree}. It operates globally over the lower ROI of the segmentation mask without strip subdivision.

Pixel values corresponding to the traversable path class are accumulated along each image column:
\begin{equation}
    H(x) = \sum_{y=y_0}^{H-1} M(y, x), \qquad x \in [0, W-1].
    \label{eq:c6_histogram}
\end{equation}
The raw profile $H(x)$ is smoothed using a one-dimensional moving-average filter to suppress high-frequency boundary irregularities:
\begin{equation}
    \tilde{H}(x) = \frac{1}{2k_w + 1} \sum_{j=-k_w}^{k_w} H(x + j).
    \label{eq:c6_histogram_smooth}
\end{equation}

The navigation reference is extracted as the column index corresponding to the profile maximum:
\begin{equation}
    x_{\text{target}} = \arg\max_{x} \tilde{H}(x).
    \label{eq:c6_histogram_target}
\end{equation}
This formulation yields a single horizontal setpoint $x_{\text{target}}$ rather than a continuous line model $x(y) = m y + q$. As a result, it does not estimate heading error $e_\theta$ directly from the visual observation, and its output is sensitive to the total area of the segmented corridor within the field of view.

\subsection{RANSAC-based contour baseline}
\label{subsec:c6_ransac_baseline}

The RANSAC-based contour baseline (\texttt{pf\_ransac.py}) reproduces the classical boundary-fitting methodology prevalent in agricultural vision navigation~\cite{bai2023vision,sensors2025deeplabv3}. Unlike Method~4 and the center-point tracker, it operates directly over the segmented corridor mask without horizontal strip discretization or bottom-up association logic.

Following lower ROI cropping ($y_0 = \lfloor 0.30 H \rfloor$) and morphological filtering, the algorithm extracts lateral boundary coordinates across all scanlines containing traversable path pixels. For each image row $y \in [y_0, H-1]$, the leftmost and rightmost corridor pixels define boundary points $(y, x_L(y))$ and $(y, x_R(y))$, provided the local corridor width $w(y) = x_R(y) - x_L(y) + 1$ satisfies $w(y) \ge 10$~px. This extracts two continuous point clouds, $\mathcal{P}_L$ and $\mathcal{P}_R$, containing all lateral boundary pixels along the corridor.

Boundary lines are estimated independently on $\mathcal{P}_L$ and $\mathcal{P}_R$ using Random Sample Consensus (RANSAC)~\cite{fischler1981ransac}:
\begin{enumerate}
    \item Across $T_{\text{max}} = 100$ trials, two points are sampled uniformly at random from the boundary point cloud.
    \item A candidate linear model $x(y) = m^{(t)} y + q^{(t)}$ is determined from the sampled pair.
    \item Inliers are identified as points satisfying $|x_k - (m^{(t)} y_k + q^{(t)})| \le \tau_{\text{res}}$, with residual threshold $\tau_{\text{res}} = 18.0$~px.
    \item The model maximizing inlier cardinality is selected, followed by an ordinary least-squares refit over the complete inlier set.
\end{enumerate}

The guidance centerline is calculated as the unweighted average of the left and right boundary models:
\begin{equation}
    m_C = \frac{m_L + m_R}{2}, \qquad q_C = \frac{q_L + q_R}{2}.
    \label{eq:c6_ransac_centerline}
\end{equation}

By fitting dense boundary clouds directly, this approach avoids heuristic strip-tuning parameters. However, evaluating candidate models across hundreds of boundary points increases computational overhead, and stochastic sampling produces non-zero angular jitter across repeated runs on identical frames.

\section{Comparative trade-offs and structural limitations}
\label{sec:c6_tradeoffs_and_limitations}

The four geometric extraction algorithms present distinct engineering trade-offs. No single method is optimal across all operational metrics. Table~\ref{tab:c6_algorithm_tradeoffs} summarizes the structural characteristics, computational properties, and operational failure modes of each extractor.

\begin{table}[htbp]
    \centering
    \footnotesize
    \setlength{\tabcolsep}{3pt}
    \renewcommand{\arraystretch}{1.2}
    \caption{Structural comparison of the four geometric extraction algorithms.}
    \label{tab:c6_algorithm_tradeoffs}
    \begin{tabularx}{\textwidth}{@{}Xcccc@{}}
        \toprule
        Characteristic & \texttt{pf\_histogram} & \texttt{pf\_centri} & \texttt{pf\_ransac} & Method~4 \\
        \midrule
        Output reference & Lateral column $x_{\text{target}}$ & Line $(m_C, q_C)$ & Line $(m_C, q_C)$ & Line $(m_C, q_C)$ \\
        Orientation ($e_\theta$) estimate & No & Yes & Yes & Yes \\
        Deterministic execution & Yes ($\sigma_\omega = 0$) & Yes ($\sigma_\omega = 0$) & No ($\sigma_\omega > 0$) & Yes ($\sigma_\omega = 0$) \\
        Computational complexity & $O(W \cdot H)$ & $O(N)$ & $O(T_{\text{max}} \cdot |\mathcal{P}|)$ & $O(N)$ \\
        Edge decoupling & N/A & No & Yes & Yes \\
        Sensitivity to asymmetric noise & High & High & Low & Low \\
        Bottom-up seed dependence & None & High & None & High \\
        Local jump threshold tuning & None & Required & None & Required \\
        Curvature representation & Poor & Moderate & Poor & Poor \\
        \bottomrule
    \end{tabularx}
\end{table}

\subsection{Structural limitations of Method~4}
\label{subsec:c6_method4_limitations}

While Method~4 decouples boundary estimation and avoids the stochastic jitter of RANSAC, it introduces specific structural constraints:

\begin{enumerate}
    \item \textbf{Bottom-up seed dependence:} The tracking process initializes from candidate intervals detected within the bottom six strips ($y \in [H-1, H - 6h_s]$). If the foreground mask immediately ahead of the vehicle contains heavy occlusions, soil debris, or severe false-negative patches, the extractor fails to identify a valid initial seed. Furthermore, if the algorithm locks onto a false boundary candidate at the base of the image, the spatial continuity constraint propagates this error upward through subsequent strips. In contrast, \texttt{pf\_histogram} and \texttt{pf\_ransac} aggregate pixels across the entire field of view simultaneously and do not depend on bottom-up temporal or spatial seeding.

    \item \textbf{Sensitivity to tracking hyperparameters:} The spatial continuity gate $\Delta x_{\text{max}} = 28$~px and the maximum allowable gap of three consecutive empty strips are tuned for a specific image resolution ($224 \times 224$) and camera mounting angle. Pitch and roll disturbances induced by uneven terrain cause rapid shifts of the corridor in image space. When these shifts exceed the spatial gate between consecutive strips, tracking terminates prematurely, resulting in frame rejection.

    \item \textbf{Linear corridor model constraint:} Method~4 models both boundaries as straight lines ($x = m y + q$). In curved crop rows, contour-following trajectories, or headland turning areas, the linear model introduces systematic approximation error. While algorithms like the multi-ROI formulation of Li et al.~\cite{li2022multiroi} adopt cubic B-splines, Method~4 trades curve-fitting capacity for numerical stability and determinism on straight corridors.

    \item \textbf{Vulnerability to extended unilateral gaps:} When crop gaps along one row exceed the three-strip threshold ($3 \times 6 = 18$~px in image space), the corresponding boundary tracker terminates. If only one boundary converges while the other is rejected, Method~4 rejects the frame unless an asymmetric single-edge fallback mode is explicitly enabled. Under identical conditions, the global histogram in \texttt{pf\_histogram} continues to output a lateral setpoint, albeit biased toward the denser vegetation side.
\end{enumerate}